% Bagian: turing-machines
% Bab: undecidability
% Subbagian: enumerating-tms

\documentclass[../../../include/open-logic-section]{subfiles}

\begin{document}

\olfileid[id]{tur}{und}{enu}
\olsection{Mengenumerasi Mesin Turing}

\begin{explain}
Kita dapat menunjukkan bahwa himpunan semua mesin Turing bersifat
!!{enumerable}. Hal ini mengikuti kenyataan bahwa setiap mesin Turing dapat
dideskripsikan secara berhingga. Himpunan keadaan dan kosakata pita merupakan
himpunan berhingga. Fungsi transisi merupakan fungsi parsial dari $Q \times
\Sigma$ ke $Q \times \Sigma \times \{\TMleft, \TMright, \TMstay\}$, sehingga
fungsi tersebut juga dapat ditentukan dengan mencantumkan nilainya bagi
sejumlah berhingga pasangan argumen tempat fungsi itu terdefinisi.

Pernyataan ini benar sejauh yang telah dikemukakan, tetapi ada suatu perbedaan
halus. Definisi mesin Turing tidak membatasi !!{element}s apa saja yang dapat
dimiliki himpunan keadaan dan alfabet pita. Jadi, misalnya, untuk setiap
bilangan real secara teknis ada sebuah mesin Turing yang menggunakan bilangan
tersebut sebagai suatu keadaan. Akan tetapi, \emph{perilaku} mesin Turing tidak
bergantung pada objek mana yang berperan sebagai keadaan dan kosakata.
Perhatikan dua mesin Turing dalam \olref{fig:variants}.
\begin{figure}
\begin{center}
\begin{tikzpicture}[->,>=stealth',shorten >=1pt,auto,node distance=2.8cm,
                    semithick]
  \tikzstyle{every state}=[fill=none,draw=black,text=black]

  \node[initial,state]         (A)              {$q_0$};
  \node[state]         (B) [right of=A] {$q_1$};

  \path (A) edge [bend left] node {\TMtrans{\TMstroke}{\TMstroke}{\TMright}} (B)
        (B) edge [loop above] node {\TMtrans{\TMblank}{\TMblank}{\TMright}} (B)
            edge [bend left] node {\TMtrans{\TMstroke}{\TMstroke}{\TMright}} (A);
\end{tikzpicture}\\
\begin{tikzpicture}[->,>=stealth',shorten >=1pt,auto,node distance=2.8cm,
  semithick]
\tikzstyle{every state}=[fill=none,draw=black,text=black]

\node[initial,state]         (A)              {$s$};
\node[state]         (B) [right of=A] {$h$};

\path (A) edge [bend left] node {\TMtrans{A}{A}{\TMright}} (B)
(B) edge [loop above] node {\TMtrans{\TMblank}{\TMblank}{\TMright}} (B)
edge [bend left] node {\TMtrans{A}{A}{\TMright}} (A);
\end{tikzpicture}
\end{center}
\caption{Varian mesin \emph{Genap}}
\ollabel{fig:variants}
\end{figure}
Kedua diagram ini bersesuaian dengan dua mesin: $M$ dengan alfabet pita
$\Sigma = \{\TMendtape,\TMblank,\TMstroke\}$ dan himpunan keadaan
$\{q_0,q_1\}$, serta $M'$ dengan alfabet $\Sigma' =
\{\TMendtape,\TMblank,A\}$ dan keadaan $\{s,h\}$. Akan tetapi, instruksi
keduanya selain itu sama: $M$ akan berhenti pada suatu urutan yang terdiri atas
$n$ simbol~$\TMstroke$ jika dan hanya jika $n$ genap, dan $M'$ akan berhenti
pada suatu urutan yang terdiri atas $n$ simbol~$A$ jika dan hanya jika $n$
genap. Yang kita lakukan hanyalah mengganti nama $\TMstroke$ menjadi~$A$,
$q_0$ menjadi~$s$, dan $q_1$ menjadi~$h$. Contoh ini dapat diperumum: kita
dapat menganggap dua mesin Turing sama selama salah satunya diperoleh dari
yang lain melalui penggantian nama simbol dan keadaan semacam itu. Bahkan,
kita dapat memandang simbol dan keadaan suatu mesin Turing sebagai bilangan
bulat positif: sebagai ganti $\sigma_0$, pikirkan~$1$; sebagai ganti
$\sigma_1$, pikirkan~$2$; dan seterusnya; $\TMendtape$ adalah~$1$,
$\TMblank$ adalah~$2$, dan seterusnya. Dengan cara ini, mesin \emph{Genap}
menjadi mesin yang digambarkan dalam \olref{fig:standard-even}.
\begin{figure}
\[\begin{tikzpicture}[->,>=stealth',shorten >=1pt,auto,node distance=2.8cm,
  semithick]
\tikzstyle{every state}=[fill=none,draw=black,text=black]

\node[initial,state]         (A)              {$1$};
\node[state]         (B) [right of=A] {$2$};

\path (A) edge [bend left] node {\TMtrans{3}{3}{\TMright}} (B)
(B) edge [loop above] node {\TMtrans{2}{2}{\TMright}} (B)
edge [bend left] node {\TMtrans{3}{3}{\TMright}} (A);
\end{tikzpicture}
\]
\caption{Mesin \emph{Genap} standar}
\ollabel{fig:standard-even}
\end{figure}
Kita dapat menyebut mesin Turing yang keadaan dan simbolnya berupa bilangan
bulat positif sebagai mesin \emph{standar}, lalu mulai sekarang hanya
mempertimbangkan mesin standar.\footnote{Istilah ``mesin standar'' bukan
istilah baku.}

Kita ingin menunjukkan bahwa himpunan mesin Turing bersifat !!{enumerable},
dan dengan mempertimbangkan hal-hal di atas, cukup ditunjukkan bahwa himpunan
mesin Turing standar bersifat !!{enumerable}. Andaikan kita diberi mesin
Turing standar $M = \tuple{Q, \Sigma, q_0, \delta}$. Bagaimana kita dapat
mendeskripsikannya dengan menggunakan suatu untaian berhingga bilangan bulat
positif? Mula-mula kita cantumkan banyaknya keadaan, keadaan-keadaan itu
sendiri, banyaknya simbol, simbol-simbol itu sendiri, dan keadaan awal.
(Ingatlah bahwa semuanya merupakan bilangan bulat positif karena $M$ adalah
mesin standar.) Bagaimana dengan~$\delta$? Himpunan argumen yang mungkin,
yakni pasangan $\tuple{q,\sigma}$, berhingga karena $Q$ dan~$\Sigma$
berhingga. Jadi, informasi dalam~$\delta$ hanyalah daftar berhingga semua
5-tupel $\tuple{q, \sigma, q', \sigma', d}$ dengan
$\delta(q,\sigma) = \tuple{q', \sigma', D}$, dan $d$ merupakan bilangan yang
mengodekan arah~$D$ (misalnya, $1$ untuk~$\TMleft$, $2$ untuk~$\TMright$,
dan $3$ untuk~$\TMstay$).

Dengan cara ini, setiap mesin Turing standar dapat dideskripsikan oleh suatu
daftar berhingga bilangan bulat positif, yakni sebagai suatu urutan $s_M \in
(\PosInt)^*$. Misalnya, mesin \emph{Genap} standar dikodekan oleh urutan
\[
2, \underbrace{1, 2}_Q, 3, \overbrace{1, 2, 3}^\Sigma, 1, \underbrace{1, 3, 2, 3, 2}_{\delta(1,3) = \tuple{2,3,R}},
\overbrace{2, 2, 2, 2, 2}^{\delta(2,2) = \tuple{2,2,R}},
\underbrace{2, 3, 1, 3, 2}_{\delta(2,3) = \tuple{1,3,R}}.
\]
\end{explain}

\begin{thm}
Ada fungsi dari $\Nat$ ke~$\Nat$ yang tidak dapat dihitung oleh mesin Turing.
\end{thm}

\begin{proof}
Kita mengetahui bahwa himpunan urutan berhingga bilangan bulat
positif~$(\PosInt)^*$ bersifat !!{enumerable}
(\cref{sfr:siz:zigzag:prob:posint-star}). Dari sini diperoleh bahwa himpunan
deskripsi mesin Turing standar, sebagai suatu subhimpunan dari~$(\PosInt)^*$,
juga terhitung. Setiap fungsi dari $\Nat$ ke~$\Nat$ yang dapat dihitung oleh
mesin Turing dihitung oleh suatu (bahkan banyak) mesin Turing. Dengan mengganti
nama keadaan dan simbolnya menjadi bilangan bulat positif (khususnya,
$\TMendtape$ menjadi~$1$, $\TMblank$ menjadi~$2$, dan $\TMstroke$ menjadi~$3$),
kita melihat bahwa setiap fungsi yang dapat dihitung oleh mesin Turing juga
dihitung oleh suatu mesin Turing standar. Ini berarti bahwa himpunan semua
fungsi dari $\Nat$ ke~$\Nat$ yang dapat dihitung oleh mesin Turing juga
terhitung.

Di pihak lain, himpunan semua fungsi dari $\Nat$ ke~$\Nat$ tidak bersifat
!!{enumerable} (\cref{sfr:siz:red:prob:nat-nat}). Seandainya semua fungsi dapat
dihitung oleh suatu mesin Turing, kita dapat mengenumerasi himpunan semua
fungsi dengan mencantumkan semua deskripsi mesin Turing yang menghitungnya.
Jadi, ada fungsi yang tidak dapat dihitung oleh mesin Turing.
\end{proof}

\begin{prob}
  Dapatkah Anda memikirkan suatu cara mendeskripsikan mesin Turing yang tidak
  mengharuskan keadaan dan simbol alfabet dicantumkan secara eksplisit? Anda
  boleh mendefinisikan sendiri gagasan mesin ``standar'', tetapi jelaskan
  mengapa setiap mesin Turing dapat disimulasikan oleh suatu mesin ``standar''
  dalam pengertian baru Anda.
\end{prob}

\end{document}
